\documentclass[10pt,a4paper]{article}
\usepackage[margin=1.7cm,top=1.3cm,bottom=1.2cm]{geometry}
\usepackage{booktabs,array,xcolor,amsmath}
\usepackage{tikz}
\usetikzlibrary{arrows.meta,calc,decorations.pathmorphing}
\pagestyle{empty}
\setlength{\parindent}{0pt}
\setlength{\parskip}{4pt}

\definecolor{good}{HTML}{1B7F4B}
\definecolor{bad}{HTML}{B3261E}
\definecolor{vox}{HTML}{3B6FB6}
\definecolor{cab}{HTML}{222222}
\definecolor{mute}{HTML}{6B6B6B}

\newcommand{\panel}[1]{{\small\bfseries #1}}
\newcolumntype{L}[1]{>{\raggedright\arraybackslash}p{#1}}
\newcommand{\pro}{\par\hangindent=1.1em\hangafter=1\makebox[1.1em][l]{\textcolor{good}{\textbf{+}}}}
\newcommand{\con}{\par\hangindent=1.1em\hangafter=1\makebox[1.1em][l]{\textcolor{bad}{\textbf{--}}}}

% one joint between two capsules; #1 stretch colour, #2 twist colour
\newcommand{\jointpic}[2]{%
\begin{tikzpicture}[x=1cm,y=1cm,>=Stealth]
  \draw[cab,line width=0.8pt,fill=black!8,rounded corners=3.2mm] (-3.1,-0.32) rectangle (-0.25,0.32);
  \draw[cab,line width=0.8pt,fill=black!8,rounded corners=3.2mm] (0.25,-0.32) rectangle (3.1,0.32);
  \fill[cab] (0,0) circle (0.07);
  % stretch
  \draw[<->,#1,line width=1pt] (-0.55,-0.62) -- (0.55,-0.62);
  % bend
  \draw[->,good,line width=1pt] (1.7,0.5) arc (20:70:1.5);
  % twist
  \draw[->,#2,line width=1pt] (2.55,-0.48) arc (-100:190:0.16 and 0.48);
\end{tikzpicture}}

\begin{document}

\begin{center}
{\Large\bfseries Cable models compared: Newton rod and PhysX capsule chain}\\[2pt]
{\small Thin cable ($d=4$\,mm), and a thick cable ($d=30$\,mm) where volumetric FEM also works}
\end{center}

\textbf{The difference in one line.} Both models are a chain of capsules. They differ
in what the joint between two capsules can do: the capsule chain has springs for
\emph{bending only}; the Newton rod has springs for \emph{stretch, bend and twist},
all derived from the material.

\vspace{3pt}
% ------------------------------------------------------------ joint pictures
\begin{minipage}[t]{0.49\linewidth}
\centering
\panel{PhysX capsule chain (D6 joint)}\\[4pt]
\jointpic{mute}{bad}\\[3pt]
{\footnotesize
\textcolor{mute}{stretch: locked}\quad
\textcolor{good}{bend: spring $EI/h$}\quad
\textcolor{bad}{twist: damper, no spring}}
\end{minipage}\hfill
\begin{minipage}[t]{0.49\linewidth}
\centering
\panel{Newton rod (\texttt{add\_rod}, VBD solver)}\\[4pt]
\jointpic{good}{good}\\[3pt]
{\footnotesize
\textcolor{good}{stretch: spring $EA/h$}\quad
\textcolor{good}{bend: spring $EI/h$}\quad
\textcolor{good}{twist: spring}}
\end{minipage}

\vspace{8pt}
\panel{1. What each joint carries}

\vspace{1pt}
\renewcommand{\arraystretch}{1.2}
{\small
\begin{tabular}{@{}L{2.6cm} L{4.6cm} L{4.6cm} L{4.6cm}@{}}
\toprule
 & \textbf{Capsule chain} (our script) & \textbf{Newton rod 1.2.0} (our script) & \textbf{Newton rod, current docs} \\
\midrule
Stretch (force along the cable) & locked: links are rigid, no axial elasticity
  & spring $EA/h$ & spring (stretch slot) \\
Bend (moment) & spring $EI/h$ on two axes & spring $EI/h$ & spring (bend slot) \\
Twist (torque about the axis) & stiffness $0$, damping only
  & carried by the joint's angular spring & its own spring (twist slot) \\
Shear & locked & not modelled & spring (shear slot) \\
Bending hysteresis & none & none & optional Dahl friction \\
Time integration & PhysX iterative, 1/480\,s step & VBD, implicit & VBD, implicit \\
\bottomrule
\end{tabular}}

\vspace{6pt}
\panel{2. Calculations, thin cable}
\quad{\small $L=893$\,mm, $r=2$\,mm, $E=80$\,MPa, $\nu=0.45$, 150 segments.}

\vspace{2pt}
{\small
\begin{minipage}[t]{0.50\linewidth}
\begin{tabular}{@{}l@{\quad}l@{}}
segment length & $h = L/150 = 893/150 = 5.95$\,mm \\
area & $A = \pi r^2 = 1.26\times10^{-5}$\,m$^2$ \\
second moment & $I = \pi r^4/4 = 1.26\times10^{-11}$\,m$^4$ \\
stretch rigidity & $EA = 1005$\,N \\
bend rigidity & $EI = 1.01\times10^{-3}$\,N\,m$^2$ \\
twist rigidity & $GJ = EI/(1+\nu) = 6.9\times10^{-4}$\,N\,m$^2$ \\
\end{tabular}
\end{minipage}\hfill
\begin{minipage}[t]{0.48\linewidth}
\begin{tabular}{@{}l@{\quad}l@{}}
stretch spring & $EA/h = 1.69\times10^{5}$\,N/m \\
bend spring & $EI/h = 0.169$\,N\,m/rad \\
twist spring & $GJ/h = 0.116$\,N\,m/rad \\[3pt]
capsule chain twist & $0$ \; (should be $0.116$) \\
capsule chain stretch & rigid \; (should be $1.69\times10^{5}$) \\
\end{tabular}
\end{minipage}}

\vspace{6pt}
\panel{3. Measured in our benchmark}
\quad{\small three supports, 150 segments, run \texttt{apple\_cable\_3pt\_150}, \texttt{metrics.csv}.}

\vspace{1pt}
{\small
\begin{center}
\begin{tabular}{@{}l c c c c c@{}}
\toprule
 & error vs catenary & error vs photo & length drift & settled & wall time \\
\midrule
Newton rod (point pin) & \textbf{0.9\,mm} & 75\,mm & 0.86\,\% & yes, 1.1\,s & 12.9\,s \\
Newton rod (tape support) & 21\,mm & 65\,mm & 1.06\,\% & yes, 1.0\,s & 11.7\,s \\
PhysX capsule chain & 74\,mm & \textbf{33\,mm} & 0.46\,\% & no, after 8.1\,s & 29.0\,s \\
\bottomrule
\end{tabular}
\end{center}}

\vspace{-2pt}
{\small The test is a planar hang: it exercises weight and bending, not twist. The
twist comparison above comes from the joint definitions, not from a measurement.}

\vspace{6pt}
\panel{4. Pros and cons, thin cable}

\vspace{2pt}
{\small
\begin{minipage}[t]{0.49\linewidth}\setlength{\parskip}{1.5pt}
\textbf{Newton rod}
\pro Torque and twist are transmitted elastically along the cable.
\pro Tension comes from $EA$, so axial force is physical.
\pro Stiffness set directly from $E$, $r$ and segment length.
\pro Implicit solver: stable at high stiffness, settles in about 1\,s.
\pro Matches the mathematics to 0.9\,mm.
\pro GPU-batched and differentiable (Warp); couples two-way with rigid-body solvers.
\con Through Isaac Sim 6 USD it cannot hold both ends (articulation must be a tree), so it runs standalone.
\con Holding a body also holds its angle: no true pin.
\con API still changing between versions.
\con Largest length drift here (0.86\,\%).
\end{minipage}\hfill
\begin{minipage}[t]{0.48\linewidth}\setlength{\parskip}{1.5pt}
\textbf{PhysX capsule chain}
\pro Native Isaac Sim: USD, Isaac Lab, ROS\,2, sensors, contact with robots.
\pro Both ends fixed works directly.
\pro Closest to the photographed cable (33\,mm).
\pro Smallest length drift (0.46\,\%), because links are rigid.
\con No twist spring in our setup: a torque spins the cable instead of loading it.
\con No axial elasticity: tension is a constraint force, not $EA\,\varepsilon$.
\con 150 rigid links solved iteratively: not settled after 8\,s, slowest run.
\con 74\,mm from the catenary.
\con Not differentiable.
\end{minipage}}

\clearpage
\panel{5. Thick cable, $d=30$\,mm: all three models work}
\quad{\small same length and material.}

\vspace{2pt}
{\small
\begin{minipage}[t]{0.42\linewidth}
\begin{tabular}{@{}l@{\quad}l@{}}
aspect ratio & $L/d = 893/30 = 30$ \\
stretch rigidity & $EA = 56.5$\,kN \\
bend rigidity & $EI = 3.18$\,N\,m$^2$ \\
twist rigidity & $GJ = 2.19$\,N\,m$^2$ \\
weight & $w = \rho A g = 16.6$\,N/m \\
bending length & $(EI/w)^{1/3} = 0.58$\,m \\
segments, $h=d$ & $893/30 \approx 30$ \\
FEM voxels across & $30/6.9 = 4.4$ at $N=130$ \\
FEM voxel count & $\tfrac{\pi}{4}\,4.4^2\cdot130 \approx 2000$ \\
\end{tabular}
\end{minipage}\hfill
\begin{minipage}[t]{0.57\linewidth}
\renewcommand{\arraystretch}{1.15}
\begin{tabular}{@{}L{2.5cm} c c c@{}}
\toprule
 & \textbf{Capsule} & \textbf{Newton rod} & \textbf{FEM} \\
\midrule
Bending & yes & yes & yes \\
Stretch & no & yes & yes \\
Twist / torque & no spring & yes & yes \\
Cross-section squeeze, contact pressure & no & no & yes \\
True geometry, no rescaling & yes & yes & yes \\
Unknowns to solve & 30 links & 30 links & $\approx$2000 voxels \\
\bottomrule
\end{tabular}
\end{minipage}}

\vspace{4pt}
{\small At 30\,mm the cable is no longer a thin string: bending length $0.58$\,m is
comparable to its $0.89$\,m length, so it behaves as a bent beam. Rod and capsule
models stay cheap but need segments about one diameter long, so they resolve
curvature with about 30 segments instead of 150. FEM costs more and is the only one that
shows the cross-section deforming, for example where a gripper squeezes it.}

\vspace{5pt}
\fcolorbox{good}{good!7}{\parbox{\dimexpr\linewidth-2\fboxsep-2\fboxrule}{%
\textbf{Summary.} Thin cable: use the Newton rod when twist, torque or tension
matter, and the capsule chain when the cable must live inside an Isaac Sim scene;
FEM is not an option. Thick cable: all three work; add FEM when the cross-section
or contact pressure matters, otherwise the rod is the cheaper choice.}}

\vspace{4pt}
{\footnotesize \textbf{Sources.} Joint definitions and stiffness formulas:
\texttt{hang\_physx\_capsule.py}, \texttt{hang\_newton\_cable.py}, \texttt{cable\_config.py} in \texttt{methods/}.
Rod stretch, shear, bend and twist slots and Dahl friction: Newton documentation, \texttt{SolverVBD}.
Two-way coupling: NVIDIA developer blog on Newton, 16 March 2026.
FEM numbers: \texttt{docs/fem\_thin\_objects}.}

\end{document}
